\documentclass[a4paper]{arrowhead}

\usepackage[yyyymmdd]{datetime}
\usepackage{etoolbox}
\usepackage[utf8]{inputenc}
\usepackage[T1]{fontenc}
\usepackage{multirow}
\usepackage{longtable}
\usepackage{tikz}
\usetikzlibrary{positioning,shapes.multipart,arrows.meta}

\renewcommand{\dateseparator}{-}
\setlength{\emergencystretch}{3em}
\lstset{language=HTTP,breakatwhitespace=false,columns=fullflexible}

%% Special references
\newcommand{\mref}[2][]{{\textcolor{ArrowheadPurple}{\hyperref[sec:model:#2]{\ifstrempty{#1}{#2}{#1}}}}}
\newcommand{\oref}[1]{{\textcolor{ArrowheadBlue}{\hyperref[sec:operations:#1]{#1}}}}
\newcommand{\pdef}[1]{{\textcolor{ArrowheadGrey}{#1\label{sec:model:primitives:#1}\label{sec:model:primitives:#1s}\label{sec:model:primitives:#1es}}}}
\newcommand{\pref}[1]{{\textcolor{ArrowheadGrey}{\hyperref[sec:model:primitives:#1]{#1}}}}

% as template_sd.tex; the input and output arguments are passed already formatted (\mref or \pref)
\newrobustcmd\osubsection[3]{
  \addtocounter{subsection}{1}
  \addcontentsline{toc}{subsection}{\protect\numberline{\thesubsection}operation \textcolor{ArrowheadBlue}{#1}}
  \subsection*{
    \thesubsection\quad
    operation
    \textcolor{ArrowheadBlue}{#1}
    (\notblank{#2}{#2}{})
    \notblank{#3}{: #3}{}
  }
  \label{sec:operations:#1}
}
\newrobustcmd\msubsection[2]{
  \addtocounter{subsection}{1}
  \addcontentsline{toc}{subsection}{\protect\numberline{\thesubsection}#1 \textcolor{ArrowheadPurple}{#2}}
  \subsection*{\thesubsection\quad#1 \textcolor{ArrowheadPurple}{#2}}
  \label{sec:model:#2} \label{sec:model:#2s} \label{sec:model:#2es}
}
%%

\usepackage{xurl}
\begin{document}
\sloppy

%% Arrowhead Document Properties
\ArrowheadTitle{translationBridge}
\ArrowheadServiceID{translationBridge}
\ArrowheadType{Service Description}
\ArrowheadTypeShort{SD}
\ArrowheadVersion{5.3.0}
\ArrowheadDate{2026-11-09}
\ArrowheadAuthor{Jerker Delsing}
\ArrowheadStatus{REVIEW DRAFT}
\ArrowheadContact{jerker.delsing@ltu.se}
\ArrowheadFooter{\href{www.arrowhead.eu}{www.arrowhead.eu}}
\ArrowheadSetup
%%

%% Front Page
\begin{center}
  \vspace*{1cm}
  \huge{\arrowtitle}

  \vspace*{0.2cm}
  \LARGE{\arrowtype}
  \vspace*{1cm}

  \Large{Service ID: \textit{"\arrowid"}}

  \vspace*{\fill}

  % Front Page Image

  \vspace*{1cm}
  \vspace*{\fill}

  % Front Page Abstract
  \begin{abstract}
    This document describes the \texttt{translationBridge} service, version 1.0.0, produced by the TranslationManager support system of Arrowhead (documentation edition 5.3.0).
  \end{abstract}

  \vspace*{1cm}
\end{center}

\newpage
%%

%% Table of Contents
\tableofcontents
\newpage
%%

\section{Documentation edition and implementation scope}
This is the 5.3.0 documentation edition, dated 2026-11-09. It carries forward the system and service contracts from the retained 5.2.2 documentation and identifies the current primary implementation sources. The former \path|documentation/v5.2.2| folder is retained unchanged.

Document version 5.3.0 is distinct from executable, dependency, registered system/service and wire-contract versions. Values still defined by the source (including 5.2.2 artefact versions, 1.0.0 service versions and \path|SysMonitor-AMQP-IDD-5.2.2| payload identifiers) are preserved. This documentation update does not change runtime behaviour or establish a new implementation release. Historical verification statements describe the tests recorded in the inherited baseline; this revision validates document compilation and references.

\section{Overview}
\label{sec:overview}

This document describes the \texttt{translationBridge} service, which provides interface and data-model translation for a consumer: finding providers that become usable through a translation bridge, building the bridge and aborting it.
It is produced by the TranslationManager system \cite{tm-sysd} in service version 1.0.0.

The rest of this document is organized as follows.
In Section \ref{sec:operations}, we describe the abstract message operations provided by the service.
In Section \ref{sec:model}, we end the document by presenting the data types used by the mentioned operations.
The concrete interfaces are specified in the HTTP/TLS/JSON IDD \cite{tm-translationBridge-http-idd} and the MQTT/TLS/JSON IDD \cite{tm-translationBridge-mqtt-idd}.

\subsection{Significant Prior Art}

Successor of the Arrowhead 4.x Translator and Gateway concepts, redesigned for Arrowhead 5 interface templates and data-model identifiers.

\subsection{How This Service Is Meant to Be Used}

Dynamic\allowbreak{}Service\allowbreak{}Orchestration normally drives translation (orchestration flag \texttt{ALLOW\_\allowbreak{}TRANSLATION}, \texttt{enable.\allowbreak{}translation=\allowbreak{}true}): \texttt{discovery} with the candidates and the consumer's interface and data models, then \texttt{negotiation} for the chosen candidate. The consumer calls the returned bridge interface instead of the provider. A consumer may also use the service directly.

\subsection{Important Delimitations and Dependencies}
\label{sec:delimitations}

Translation needs interface and data-model translator systems registered in the cloud; the Translation\allowbreak{}Manager orchestrates them but does not translate itself.

\newpage

\section{Service Interface}
\label{sec:operations}

This section describes the interfaces to the \texttt{translationBridge} service.
In particular, each subsection names an abstract operation, an input type and an output type, in that order.
The input type is named inside parentheses, while the output type is preceded by a colon.
Input and output types are only denoted when accepted or returned, respectively, by the interface in question.

All abstract data types named in this section are defined in Section \ref{sec:model}.

\begin{figure}[ht!]
  \centering
  \begin{tikzpicture}[font=\small]
    \node[draw=ArrowheadBlue, very thick, rectangle split, rectangle split parts=2, align=center,
          rectangle split part fill={ArrowheadBlue!12,white}, text width=9cm, inner sep=6pt] (b)
      {\guillemotleft{}Service\guillemotright{}\\\textbf{translationBridge}
        \nodepart{second} \texttt{discovery}\\ \texttt{negotiation}\\ \texttt{abort}};
  \end{tikzpicture}
  \caption{SysML-style block of the \texttt{translationBridge} service and its operations.}
  \label{fig:SD}
\end{figure}

The following interface operations are available.

\osubsection{discovery}{\mref[Translation\allowbreak{}Discovery\allowbreak{}Request]{TranslationDiscoveryRequest}}{\mref[Translation\allowbreak{}Discovery\allowbreak{}Response]{TranslationDiscoveryResponse}}

Returns the candidates usable by the requester through translation, with a bridge identifier.

\osubsection{negotiation}{\mref[Translation\allowbreak{}Negotiation\allowbreak{}Request]{TranslationNegotiationRequest}}{\mref[Translation\allowbreak{}Negotiation\allowbreak{}Response]{TranslationNegotiationResponse}}

Builds the bridge to the chosen candidate and returns the bridge interface and token limits.

\osubsection{abort}{\pref{Identifier}}{}

Aborts a bridge of the requester. Returns 200 if aborted, 204 if there was nothing to abort. The bridge identifier fills the \texttt{\{bridgeId\}} placeholder of the registered path.

\newpage

\section{Information Model}
\label{sec:model}

Here, all data objects that can be part of the \texttt{translationBridge} service are listed in alphabetic order.
Note that each subsection, which describes one type of object, begins with the \textit{struct} keyword, which is used to denote a collection of named fields, each with its own data type.
As a complement to the explicitly defined types in this section, there is also a list of implicit primitive types in Section \ref{sec:model:primitives}, which are used to represent things like hashes and identifiers.

\msubsection{struct}{Address}

A network address with its type.


\noindent\begin{tabularx}{\textwidth}{| >{\raggedright\arraybackslash}p{3.9cm} | >{\raggedright\arraybackslash\small}p{5.3cm} | X |} \hline
\rowcolor{gray!33} Field & Type & Description \\ \hline
type & \pref{String} & Address type: \texttt{HOSTNAME}, \texttt{IPV4}, \texttt{IPV6} or \texttt{MAC}. \\ \hline
address & \pref{String} & One network address (host name, IPv4, IPv6 or MAC address). \\ \hline
\end{tabularx}



\msubsection{struct}{DeviceResponse}

A registered device.


\noindent\begin{tabularx}{\textwidth}{| >{\raggedright\arraybackslash}p{3.9cm} | >{\raggedright\arraybackslash\small}p{5.3cm} | X |} \hline
\rowcolor{gray!33} Field & Type & Description \\ \hline
name & \pref{Name} & Device name. \\ \hline
metadata & \pref{Metadata} & Free-form metadata (JSON object). \\ \hline
addresses & \pref{List}$<$\allowbreak{}\mref{Address}$>$ & Network addresses with their types. \\ \hline
created\allowbreak{}At & \pref{DateTime} & UTC time of creation. \\ \hline
updated\allowbreak{}At & \pref{DateTime} & UTC time of the last change. \\ \hline
\end{tabularx}



\msubsection{struct}{ServiceDefinitionResponse}

A service definition.


\noindent\begin{tabularx}{\textwidth}{| >{\raggedright\arraybackslash}p{3.9cm} | >{\raggedright\arraybackslash\small}p{5.3cm} | X |} \hline
\rowcolor{gray!33} Field & Type & Description \\ \hline
name & \pref{Name} & Service definition name. \\ \hline
created\allowbreak{}At & \pref{DateTime} & UTC time of creation. \\ \hline
updated\allowbreak{}At & \pref{DateTime} & UTC time of the last change. \\ \hline
\end{tabularx}



\msubsection{struct}{ServiceInstanceInterfaceResponse}

One interface of a service instance.


\noindent\begin{tabularx}{\textwidth}{| >{\raggedright\arraybackslash}p{3.9cm} | >{\raggedright\arraybackslash\small}p{5.3cm} | X |} \hline
\rowcolor{gray!33} Field & Type & Description \\ \hline
template\allowbreak{}Name & \pref{Name} & Name of the interface template, snake\_case (e.g. \texttt{generic\_\allowbreak{}https}). \\ \hline
protocol & \pref{String} & Protocol of the interface, e.g. \texttt{http}, \texttt{https}, \texttt{mqtt}, \texttt{mqtts}. \\ \hline
policy & \pref{String} & Security policy of the interface: \texttt{NONE}, \texttt{CERT\_\allowbreak{}AUTH} or one of the token policies. \\ \hline
properties & \pref{Map}$<$\allowbreak{}\pref{String}, \pref{Object}$>$ & Interface properties defined by the template, e.g. \texttt{accessAddresses}, \texttt{accessPort}, \texttt{basePath}, \texttt{operations}. \\ \hline
\end{tabularx}



\msubsection{struct}{ServiceInstanceResponse}

A registered service instance.


\noindent\begin{tabularx}{\textwidth}{| >{\raggedright\arraybackslash}p{3.9cm} | >{\raggedright\arraybackslash\small}p{5.3cm} | X |} \hline
\rowcolor{gray!33} Field & Type & Description \\ \hline
instance\allowbreak{}Id & \pref{Identifier} & Service instance identifier \texttt{\textless{}System\textgreater{}\textbar{}\allowbreak{}\textless{}service\allowbreak{}Definition\textgreater{}\textbar{}\allowbreak{}\textless{}version\textgreater{}}, or policy instance identifier. \\ \hline
provider & \mref{SystemResponse} & The provider system. \\ \hline
service\allowbreak{}Definition & \mref[Service\allowbreak{}Definition\allowbreak{}Response]{ServiceDefinitionResponse} & The service definition. \\ \hline
version & \pref{Version} & Version of the service instance. \\ \hline
expires\allowbreak{}At & \pref{DateTime} & UTC expiry time; no value means no expiry. \\ \hline
metadata & \pref{Metadata} & Free-form metadata (JSON object). \\ \hline
interfaces & \pref{List}$<$\allowbreak{}\mref[Service\allowbreak{}Instance\allowbreak{}Interface\allowbreak{}Response]{ServiceInstanceInterfaceResponse}$>$ & Service interfaces offered by the instance. \\ \hline
created\allowbreak{}At & \pref{DateTime} & UTC time of creation. \\ \hline
updated\allowbreak{}At & \pref{DateTime} & UTC time of the last change. \\ \hline
\end{tabularx}



\msubsection{struct}{SystemResponse}

A registered system.


\noindent\begin{tabularx}{\textwidth}{| >{\raggedright\arraybackslash}p{3.9cm} | >{\raggedright\arraybackslash\small}p{5.3cm} | X |} \hline
\rowcolor{gray!33} Field & Type & Description \\ \hline
name & \pref{Name} & System name. \\ \hline
metadata & \pref{Metadata} & Free-form metadata (JSON object). \\ \hline
version & \pref{Version} & Version of the system. \\ \hline
addresses & \pref{List}$<$\allowbreak{}\mref{Address}$>$ & Network addresses with their types. \\ \hline
device & \mref{DeviceResponse} & The device hosting the system, if registered. \\ \hline
created\allowbreak{}At & \pref{DateTime} & UTC time of creation. \\ \hline
updated\allowbreak{}At & \pref{DateTime} & UTC time of the last change. \\ \hline
\end{tabularx}



\msubsection{struct}{TranslationBridgeCandidate}

A candidate reachable through translation.


\noindent\begin{tabularx}{\textwidth}{| >{\raggedright\arraybackslash}p{3.9cm} | >{\raggedright\arraybackslash\small}p{5.3cm} | X |} \hline
\rowcolor{gray!33} Field & Type & Description \\ \hline
service\allowbreak{}Instance\allowbreak{}Id & \pref{Identifier} & Service instance identifier \texttt{\textless{}System\textgreater{}\textbar{}\allowbreak{}\textless{}service\allowbreak{}Definition\textgreater{}\textbar{}\allowbreak{}\textless{}version\textgreater{}}. \\ \hline
interface\allowbreak{}Template\allowbreak{}Name & \pref{String} & Interface template of the candidate the bridge would translate to. \\ \hline
\end{tabularx}



\msubsection{struct}{TranslationDiscoveryRequest}

Discovery of translation bridges for the requester.


\noindent\begin{tabularx}{\textwidth}{| >{\raggedright\arraybackslash}p{3.9cm} | >{\raggedright\arraybackslash\small}p{5.3cm} | X |} \hline
\rowcolor{gray!33} Field & Type & Description \\ \hline
candidates & \pref{List}$<$\allowbreak{}\mref{ServiceInstanceResponse}$>$ & Candidate service instances (from service discovery). \\ \hline
operation & \pref{String} & Operation the consumer wants to call. \\ \hline
interface\allowbreak{}Template\allowbreak{}Names & \pref{List}$<$\allowbreak{}\pref{String}$>$ & Interface templates the consumer can use. \\ \hline
input\allowbreak{}Data\allowbreak{}Model\allowbreak{}Id & \pref{Identifier} & Data model of the operation input expected by the consumer. \\ \hline
output\allowbreak{}Data\allowbreak{}Model\allowbreak{}Id & \pref{Identifier} & Data model of the operation output expected by the consumer. \\ \hline
\end{tabularx}



\msubsection{struct}{TranslationDiscoveryResponse}

Candidates reachable through a translation bridge.


\noindent\begin{tabularx}{\textwidth}{| >{\raggedright\arraybackslash}p{3.9cm} | >{\raggedright\arraybackslash\small}p{5.3cm} | X |} \hline
\rowcolor{gray!33} Field & Type & Description \\ \hline
bridge\allowbreak{}Id & \pref{Identifier} & Identifier (UUID) of the translation bridge. \\ \hline
candidates & \pref{List}$<$\allowbreak{}\mref[Translation\allowbreak{}Bridge\allowbreak{}Candidate]{TranslationBridgeCandidate}$>$ & Store entries to create. \\ \hline
\end{tabularx}



\msubsection{struct}{TranslationNegotiationRequest}

Builds the bridge to the chosen candidate.


\noindent\begin{tabularx}{\textwidth}{| >{\raggedright\arraybackslash}p{3.9cm} | >{\raggedright\arraybackslash\small}p{5.3cm} | X |} \hline
\rowcolor{gray!33} Field & Type & Description \\ \hline
bridge\allowbreak{}Id & \pref{Identifier} & Identifier (UUID) of the translation bridge. \\ \hline
target & \mref{ServiceInstanceResponse} & The chosen candidate service instance. \\ \hline
operation & \pref{String} & QoS operation: \texttt{FILTER} or \texttt{SORT}. \\ \hline
interface\allowbreak{}Template\allowbreak{}Name & \pref{String} & Interface template the consumer will use. \\ \hline
input\allowbreak{}Data\allowbreak{}Model\allowbreak{}Id & \pref{Identifier} & Data model of the operation input expected by the consumer. \\ \hline
output\allowbreak{}Data\allowbreak{}Model\allowbreak{}Id & \pref{Identifier} & Data model of the operation output expected by the consumer. \\ \hline
\end{tabularx}



\msubsection{struct}{TranslationNegotiationResponse}

The established bridge.


\noindent\begin{tabularx}{\textwidth}{| >{\raggedright\arraybackslash}p{3.9cm} | >{\raggedright\arraybackslash\small}p{5.3cm} | X |} \hline
\rowcolor{gray!33} Field & Type & Description \\ \hline
bridge\allowbreak{}Id & \pref{Identifier} & Identifier (UUID) of the translation bridge. \\ \hline
bridge\allowbreak{}Interface & \mref[Service\allowbreak{}Instance\allowbreak{}Interface\allowbreak{}Response]{ServiceInstanceInterfaceResponse} & Interface of the bridge the consumer calls instead of the provider. \\ \hline
token\allowbreak{}Expires\allowbreak{}At & \pref{DateTime} & Expiry of the bridge token. \\ \hline
token\allowbreak{}Usage\allowbreak{}Limit & \pref{Number} & Number of uses allowed with the bridge token. \\ \hline
\end{tabularx}



\subsection{Primitives}
\label{sec:model:primitives}

Types and structures mentioned throughout this document that are assumed to be available to implementations of this service.
The concrete interpretations of each of these types and structures must be provided by any IDD document claiming to implement this service.

\noindent\begin{tabularx}{\textwidth}{| p{4.6cm} | X |} \hline
\rowcolor{gray!33} Type & Description \\ \hline
\pdef{DateTime} & Pinpoints a specific moment in time (UTC). \\ \hline
\pdef{Identifier} & A string that identifies a record, e.g. a UUID or a composite identifier. \\ \hline
\pdef{List} & An \textit{array} of a known number of items, each having type A. \\ \hline
\pdef{Map} & A set of key-value pairs. \\ \hline
\pdef{Metadata} & A set of key-value pairs with arbitrary, possibly nested values. \\ \hline
\pdef{Name} & A string identifier that is intended to be both human and machine-readable. \\ \hline
\pdef{Number} & An integer number. \\ \hline
\pdef{Object} & Any structured value. \\ \hline
\pdef{String} & A chain of characters. \\ \hline
\pdef{Version} & Specifies a version, semantic versioning \texttt{major.minor.patch}. \\ \hline
\end{tabularx}

\newpage

\bibliographystyle{IEEEtran}
\bibliography{bibliography}

\newpage

\section{Revision History}
\subsection{Amendments}

\noindent\begin{tabularx}{\textwidth}{| p{1cm} | p{3cm} | p{2cm} | X | p{4cm} |} \hline
\rowcolor{gray!33} No. & Date & Version & Subject of Amendments & Author \\ \hline

1 & 2026-10-22 & 5.2.2 & Document created from the Arrowhead LaTeX templates for the 5.2.2 Java Spring implementation, including the local extensions of this workspace. & Jerker Delsing \\ \hline
2 & 2026-10-30 & 5.2.2 & Checked against the primary Java source tree (Eclipse-Arrowhead-java): services, operations, paths and data models unchanged. & Claude Opus 5.5 (draft) \\ \hline
3 & 2026-11-09 & 5.3.0 & New documentation edition; primary-source references and Go/Rust implementation scope; 5.2.2 retained. & Documentation maintenance \\ \hline
\end{tabularx}

\subsection{Quality Assurance}

\noindent\begin{tabularx}{\textwidth}{| p{1cm} | p{3cm} | p{2cm} | X |} \hline
\rowcolor{gray!33} No. & Date & Version & Approved by \\ \hline

1 & & \arrowversion  &  \\ \hline

\end{tabularx}

\end{document}
